%================================
% Vocabulary & Language Focus — XB3 (Unit 4)
% 编译方式: xelatex main.tex (需运行两遍以生成目录)
%================================
\documentclass[../main.tex]{subfiles}

\begin{document}

\section{Unit 4}

\vw{heritage}{/ˈherɪtɪdʒ/}{n.}

\sense{n.}{遗产}
\col{cultural/natural heritage 文化遗产/自然遗产}
\ex{The city is famous for its cultural heritage, including ancient temples and medieval castles.}
\col{world heritage site 世界遗产}
\ex{The Great Wall is a UNESCO World Heritage Site.}
\col{a heritage of sth}
\ex{The two countries share a common heritage of language and culture.}

\begin{origin}
13 世纪借自古法语 \textit{heritage}（继承之物），源于 \textit{heriter}（继承），拉丁语 \textit{hereditare}（继承、传袭），最终追溯至拉丁语 \textit{heres}（继承人）。同源词 \textbf{inherit}（继承；\textit{in-} 入 + \textit{hereditare} $\to$ 进入继承$\to$ 接受遗产）、\textbf{inheritance}（遗产；继承）、\textbf{heir}（继承人）、\textbf{heirloom}（传家宝；\textit{heir} 继承人 + \textit{loom} 工具——继承的工具）。"heritage" 字面"被继承之物"——文化、传统、家族财产皆可。
\end{origin}

\vw{awe}{/ɔː/}{vt. \& n.}

\sense{vt.}{使惊叹，使敬畏}
\col{awe sb}
\ex{The magnificent view from the summit awed every member of the climbing team.}
\col{be awed by/into sth}
\ex{The children were awed into silence by the vastness of the cathedral.}

\sense{n.}{敬畏，惊叹}
\col{in awe of sth}
\ex{We stood in awe of the towering cliffs before us.}
\col{be/stand in awe of sb/sth}
\ex{She has always stood in awe of her grandmother's strength of character.}

\begin{gram}{awe / fear / wonder / reverence 的"敬畏"细分}
\textbf{awe}（敬畏——含"惊叹"和"畏惧"双重情感）；\textbf{fear}（恐惧——纯负面）；\textbf{wonder}（惊叹——偏正面好奇）；\textbf{reverence}（敬畏——偏宗教/道德崇敬）。awe 是敬畏的"中性"层面——既惊又惧。
\end{gram}

\begin{origin}
13 世纪中古英语 \textit{age}（恐怖、敬畏），古英语 \textit{ege}（恐惧），原始日耳曼语 *\textit{agiz}。同源词 \textbf{awe}、\textbf{awesome}（令人敬畏的；极好的）、\textbf{awful}（糟糕的；古义"令人敬畏的"）、\textbf{awkward}（尴尬的；古义"方向错误"）。awful 与 awesome 同根——古英语中"令人敬畏"是中性的褒贬兼有。
\end{origin}

\vw{civilization}{/ˌsɪvəlaɪˈzeɪʃn/}{n.}

\sense{n.}{文明；社会文明}
\col{ancient/modern/Western civilization}
\ex{The ancient civilizations of Egypt, Greece and China have shaped the modern world.}
\col{the cradle of civilization 文明摇篮}
\ex{The Fertile Crescent is often called the cradle of civilization.}

\sense{n.}{文明社会，文明世界}
\col{return to civilization（幽默用法）}
\ex{After three weeks in the wilderness, it was a relief to return to civilization.}

\begin{origin}
18 世纪由 \textit{civilize}（使开化）+ \textit{-ation} 派生。civilize 17 世纪借自法语 \textit{civiliser}，源于拉丁语 \textit{civilis}（属于公民的；文明的），是 \textit{civis}（公民）的形容词形式。同源词 \textbf{civil}（公民的；民事的；文明的）、\textbf{civic}（公民的；城市的）、\textbf{citizen}（公民）、\textbf{city}（城市）、\textbf{uncivilized}（未开化的；野蛮的）、\textbf{uncivil}（不礼貌的；\textit{un-} 不 + \textit{civilis} $\to$ 不像公民的）。
\end{origin}

\vw{preserve}{/prɪˈzɜːv/}{vt.}

\sense{vt.}{维护，保护，保存}
\col{preserve sth}
\ex{We must do more to preserve the world's endangered species.}
\col{preserve the peace/order}
\ex{The army was called in to preserve order during the demonstration.}

\sense{vt.}{维持}
\col{preserve one's health/independence/way of life}
\ex{They moved to the countryside to preserve their traditional way of life.}

\begin{lex}{preserve / conserve / reserve / save 的"保"细分}
\textbf{preserve}（保存——使不变质：\textit{preserve food in vinegar} 用醋保存食物）；\textbf{conserve}（保护——节约使用：\textit{conserve energy} 节约能源）；\textbf{reserve}（保留；预订：\textit{reserve a seat}）；\textbf{save}（拯救——多指危难：\textit{save lives} 救命）。preserve 偏"使维持原样"，conserve 偏"节约资源"。
\end{lex}

\begin{origin}
14 世纪借自拉丁语 \textit{praeservare}（守卫、保存）：\textit{prae-}（在前）+ \textit{servare}（保护、保存）。同源词 \textbf{preserve}、\textbf{preservation}（保护）、\textbf{preservative}（防腐剂）、\textbf{observe}（观察；遵守；\textit{ob-} 朝 + servare $\to$ 面向保护$\to$ 观察/遵守）、\textbf{conserve}（保存）、\textbf{conserve}（节约）、\textbf{reserve}（保留）、\textbf{deserve}（应得；\textit{de-} 加强 + servare $\to$ 充分服务$\to$应得）、\textbf{servant}（仆人；\textit{servire} 服务）。
\end{origin}

\vw{urgent}{/ˈɜːdʒənt/}{adj.}

\sense{adj.}{紧急的，紧迫的}
\col{an urgent matter/message/case}
\ex{We have an urgent matter to discuss before the meeting ends.}
\col{be in urgent need of sth}
\ex{The hospital is in urgent need of blood donations.}

\sense{adj.}{催促的，急切的}
\col{an urgent tone of voice}
\ex{There was an urgent tone in her voice that made me worry.}

\begin{gram}{urgent / pressing / critical / acute 的"急"细分}
\textbf{urgent}（紧迫的——需立即行动：\textit{an urgent message}）；\textbf{pressing}（迫切的——需尽快：\textit{a pressing problem}）；\textbf{critical}（危急的——多指生死攸关：\textit{a critical condition} 危急情况）；\textbf{acute}（剧烈的；急性的——多指疾病：\textit{an acute pain}）。考试常考：critical condition vs. acute pain——病情的不同角度。
\end{gram}

\begin{origin}
15 世纪借自拉丁语 \textit{urgentem}（驱迫的、急切的），是 \textit{urgere}（驱迫、强迫）的现在分词。同源词 \textbf{urge}（催促；强烈欲望）、\textbf{urgency}（紧急）、\textbf{urgently}（紧急地）。"urgent" 字面"被驱迫的"——内心有一股力量推着你立刻行动。
\end{origin}

\vp{wear down}

\sense{短语}{(使) 磨损，(使) 逐渐磨平}
\col{wear down sth}
\ex{The river has worn down the rocks over thousands of years.}
\col{wear sb down 消磨某人的意志}
\ex{The long negotiations were designed to wear down the opposition's resistance.}

\begin{gram}{wear 短语小家族}
\textbf{wear down}（磨损；消磨）；\textbf{wear off}（逐渐消失：\textit{The effect of the drug will wear off after a few hours}）；\textbf{wear out}（穿坏；耗尽：\textit{wear out one's welcome} 不再受欢迎）；\textbf{wear away}（磨掉；消逝）；\textbf{wear thin}（变薄；减弱：\textit{My patience is wearing thin}）。
\end{gram}

\vw{ruin}{/ˈruːɪn/}{vt. \& n.}

\sense{vt.}{破坏，毁坏}
\col{ruin sth}
\ex{The earthquake ruined the city, leaving thousands homeless.}
\col{ruin sb's life/plan/holiday}
\ex{The heavy rain completely ruined our holiday.}

\sense{vt.}{使破产，毁灭}
\col{ruin sb}
\ex{Failing the exam ruined his hopes of going to university.}

\sense{n.}{毁坏，毁灭}
\col{in ruins 沦为废墟}
\ex{After the war, the city lay in ruins.}
\col{financial ruin}
\ex{She was facing financial ruin after losing her job.}

\sense{n.}{破产}
\col{be on the brink/verge of ruin}
\ex{The company is on the brink of ruin.}

\sense{n.}{残垣断壁}
\col{the ruins of ...}
\ex{We visited the ancient ruins of a Greek temple.}

\begin{gram}{ruin / destroy / demolish / wreck 的"毁"细分}
\textbf{ruin}（毁坏——多指渐进或全面：\textit{ruin one's health}）；\textbf{destroy}（摧毁——强力破坏：\textit{destroy the building}）；\textbf{demolish}（拆毁——多指建筑：\textit{demolish the old house}）；\textbf{wreck}（失事——多指车辆/船：\textit{wreck the car}）。
\end{gram}

\begin{origin}
14 世纪借自古法语 \textit{ruine}，拉丁语 \textit{ruina}（倒塌、崩溃），源于 \textit{ruere}（倒下、奔向）。同源词 \textbf{ruin}、\textbf{ruinous}（毁灭性的；\textit{ruinous expense} 巨额费用）、\textbf{ruins}（废墟）、\textbf{prone}（俯卧的；倾向——\textit{pro-} 前 + ruere $\to$ 向前倒下$\to$ 俯卧的）、\textbf{crash}（崩溃——与 ruin 形近但实不同源）。
\end{origin}

\vw{incredible}{/ɪnˈkredəbl/}{adj.}

\sense{adj.}{极好的，极大的}
\col{an incredible view/performance/experience}
\ex{The view from the top of the mountain was absolutely incredible.}
\ex{She has an incredible talent for languages.}

\sense{adj.}{不能相信的，难以置信的}
\col{an incredible story/excuse}
\ex{That excuse sounds pretty incredible to me — are you sure it's true?}

\begin{gram}{incredible / incredulous / uncredible 的细微差异}
\textbf{incredible}（难以置信的；极好的——形容词，描述事物）；\textbf{incredulous}（不相信的；表示怀疑——多形容人：\textit{an incredulous look}）；\textbf{uncredible}（不可信的——不常用）。考试常考：She gave me an incredulous look.（她投以怀疑的目光）。
\end{gram}

\begin{origin}
15 世纪借自拉丁语 \textit{incredibilis}（不可相信的）：\textit{in-}（不）+ \textit{credibilis}（可信的），源于 \textit{credere}（相信）。同源词 \textbf{credible}（可信的）、\textbf{credit}（信用；学分；\textit{credit card} 信用卡）、\textbf{creed}（信条；信经）、\textbf{credential}（凭证）、\textbf{discredit}（怀疑；损害名誉；\textit{dis-} 失去 + credit $\to$ 失去信用）。incredible 在美式口语中已"褒义化"——表"极好的"。
\end{origin}

\vw{sculpture}{/ˈskʌlptʃə(r)/}{n.}

\sense{n.}{雕像，雕塑品}
\col{a marble/bronze/wooden sculpture}
\ex{The exhibition features marble sculptures from ancient Greece.}
\col{by sculpture}
\ex{The figure is represented in sculpture rather than painting.}

\sense{n.}{雕刻术}
\col{the art of sculpture}
\ex{She studied the art of sculpture at the Royal Academy.}

\begin{origin}
17 世纪借自拉丁语 \textit{sculptura}，是 \textit{sculpere}（雕刻、刮）的派生形式。同源词 \textbf{sculpt}（雕塑；雕刻）、\textbf{sculptor}（雕塑家）、\textbf{sculptural}（雕塑的；雕刻的）、\textbf{inscription}（铭文；\textit{in-} 在上 + \textit{scribere} 写——与 sculpture 形近但实不同源，scribere 写 vs. sculpere 刻）。sculpture 与 inscription 的"刻"与"写"——古罗马人不分。
\end{origin}

\vw{missile}{/ˈmɪsaɪl/}{n.}

\sense{n.}{导弹}
\col{a guided/ballistic/nuclear missile}
\ex{The country has developed long-range ballistic missiles capable of reaching the other side of the continent.}
\col{missile attack/strike 导弹攻击}
\ex{The capital was hit by a missile strike overnight.}

\begin{origin}
17 世纪借自拉丁语 \textit{missile}（可被抛出的；投射物），是 \textit{mittere}（发送、发射）的形容词形式。同源词 \textbf{mission}（任务）、\textbf{missionary}（传教士）、\textbf{mission statement}（使命宣言）、\textbf{admit}（承认；\textit{ad-} 朝 + mittere $\to$ 朝向发送$\to$允许入内）、\textbf{emit}（发射）、\textbf{submit}（提交；屈服）、\textbf{transmit}（传输）、\textbf{dismiss}（解雇；解散）、\textbf{promise}（承诺；\textit{pro-} 前 + mittere $\to$预先放出之言）。
\end{origin}

\vw{violence}{/ˈvaɪələns/}{n.}

\sense{n.}{暴力，暴行}
\col{domestic/gang/gun violence}
\ex{Domestic violence is a serious social problem in many countries.}
\col{an act of violence 暴力行为}
\ex{The government condemned the act of violence against civilians.}
\col{resort to violence 诉诸暴力}
{When negotiations failed, the protesters resorted to violence.}

\sense{n.}{狂热，激情}
\col{with violence}
\ex{The waves dashed against the rocks with tremendous violence.}

\begin{gram}{violence / force / aggression 的"暴"细分}
\textbf{violence}（暴力——使用强力伤害：\textit{domestic violence}）；\textbf{force}（力量——中性：\textit{use force} 使用武力）；\textbf{aggression}（侵略——主动攻击：\textit{military aggression}）。violence 偏具体伤害，force 偏中性的"力"，aggression 偏敌意行为。
\end{gram}

\begin{origin}
14 世纪借自古法语 \textit{ violence}，拉丁语 \textit{violentia}（猛烈、暴力），源于 \textit{vis}（力量）。同源词 \textbf{violent}（暴力的）、\textbf{violently}（猛烈地）、\textbf{violate}（违反；侵犯）、\textbf{violinist}（小提琴手——与 violent 形近但实不同源，violin 来自中世纪拉丁语 \textit{vitula} 弦乐器）。
\end{origin}

\vw{conflict}{/ˈkɒnflɪkt/ n. \& /kənˈflɪkt/ v.}

\sense{n.}{(军事) 冲突，战斗}
\col{armed conflict 武装冲突}
\ex{The two countries have been in armed conflict for over a decade.}
\col{a military conflict}
\ex{The military conflict has caused thousands of casualties.}

\sense{n.}{争论}
\col{conflict over/about sth}
{There is ongoing conflict over the proposed new airport.}

\sense{n.}{矛盾，不一致}
\col{conflict between A and B}
\ex{There is a clear conflict between the two sides' objectives.}
\col{in conflict with sth}
{These findings are in direct conflict with earlier research.}

\sense{vi.}{冲突，抵触}
\col{conflict with sth}
{Our political views conflict with each other on this issue.}

\begin{gram}{conflict 的重音转移：名词 vs. 动词}
\textbf{conflict}（n. 冲突）重音在前 /ˈkɒnflɪkt/；\textbf{conflict}（vi. 冲突）重音在后 /kənˈflɪkt/。同类的重音转移词：\textbf{record}、\textbf{conduct}、\textbf{contrast}、\textbf{increase}、\textbf{decrease}、\textbf{project} 等。
\end{gram}

\begin{origin}
15 世纪借自拉丁语 \textit{conflictus}（撞击、冲突），是 \textit{confligere}（撞击、冲突）的过去分词：\textit{con-}（一起）+ \textit{fligere}（打击）。同源词 \textbf{conflict}、\textbf{confliction}（冲突；少见）、\textbf{afflict}（使痛苦；\textit{af-} 朝 + \textit{fligere} $\to$ 朝向打击$\to$ 折磨）、\textbf{inflict}（施加；强加；\textit{in-} 在上 + fligere $\to$ 击打上去）。"conflict" 字面"互相撞击"。
\end{origin}

\vw{civil}{/ˈsɪvl/}{adj.}

\sense{adj.}{国民的；国家的}
\col{civil rights/duties 公民权利/义务}
\ex{The civil rights movement of the 1960s changed American society forever.}
\col{civil war 内战}
\ex{The country is still recovering from a long and bloody civil war.}

\sense{adj.}{民事的}
\col{civil case/law 民事案件/民法}
\ex{The dispute was settled in a civil court rather than a criminal one.}
\col{civil servant 公务员}
\ex{She has worked as a civil servant for over twenty years.}

\begin{gram}{civil / civic / civilian / civilised 的"民"分叉}
\textbf{civil}（公民的；民事的；\textit{civil rights}）；\textbf{civic}（城市的；公民的——多形容责任/义务：\textit{civic duty}）；\textbf{civilian}（平民——非军人：\textit{civilian casualties}）；\textbf{civilised}（文明的；开化的）。civil 偏法律/权利，civic 偏责任/参与。
\end{gram}

\begin{origin}
14 世纪借自古法语 \textit{civil}，拉丁语 \textit{civilis}（属于公民的；礼貌的），源于 \textit{civis}（公民）。同源词 \textbf{civil}、\textbf{citizen}（公民）、\textbf{city}（城市）、\textbf{civic}（公民的；城市的）、\textbf{civilian}（平民）、\textbf{civilization}（文明）、\textbf{uncivil}（不礼貌的）、\textbf{civility}（礼貌）。
\end{origin}

\vw{status}{/ˈsteɪtəs/}{n.}

\sense{n.}{身份，地位}
\col{social/economic/marital status}
{Your economic status will affect your access to credit.}
\col{seek status 追求地位}
{She worked hard to achieve a higher social status for her family.}

\sense{n.}{状况，情形}
\col{status quo 现状}
{We are not happy with the status quo — we want change.}
\col{status symbol 身份象征}
{The expensive sports car is a status symbol for many wealthy people.}

\begin{gram}{status / statue / stature / state / statute 的"stat-"词族分叉}
\textbf{status}（地位；身份；现状）；\textbf{statue}（雕像——"站立之物"）；\textbf{stature}（身高；声望——"站立的高度"）；\textbf{state}（状态；国家；陈述）；\textbf{statute}（法令；法规——"已制定之物"）。五者皆源自拉丁语 \textit{stare}（站立），但引申义各异。
\end{gram}

\begin{origin}
17 世纪借自拉丁语 \textit{status}（站立、姿势、地位、状态），是 \textit{stare}（站立）的过去分词。同源词 \textbf{status}、\textbf{status quo}（现状；拉丁语"状态在何处"）、\textbf{state}（状态；国家）、\textbf{statement}（陈述）、\textbf{station}（车站；站）、\textbf{statue}（雕像）、\textbf{static}（静态的）、\textbf{statistics}（统计学）、\textbf{estate}（地产——"站立之物"）。"status" 字面"站的位置"——社会空间中的站位。
\end{origin}

\vw{internal}{/ɪnˈtɜːnl/}{adj.}

\sense{adj.}{内部的}
\col{internal organs/injuries}
{She suffered serious internal injuries in the accident.}
\col{an internal examination 内部检查}
{The committee will conduct an internal examination of the financial records.}

\sense{adj.}{国内的}
\col{internal affairs/trade}
{The country is facing serious problems on the internal front.}
\col{internal flights 国内航班}
{From Beijing, internal flights connect to all major cities in China.}

\begin{gram}{internal / interior / inner / inward / inside 的"内"细分}
\textbf{internal}（内部的——多形容抽象/组织：\textit{internal affairs}）；\textbf{interior}（内部的——多形容物理空间：\textit{the interior of the building}）；\textbf{inner}（内心的；内圈的——\textit{inner circle} 核心圈）；\textbf{inward}（向内的；内心的——\textit{inward thoughts}）；\textbf{inside}（里面的；内部的——口语化）。internal 偏组织/系统，interior 偏物理空间。
\end{gram}

\begin{origin}
15 世纪由 \textit{inter-}（内、中间）+ \textit{-al} 派生。inter- 借自拉丁语 \textit{inter}（在内、在……之间）。同源词 \textbf{interior}（内部）、\textbf{internal}（内部的）、\textbf{internalize}（使内化）、\textbf{international}（国际的；\textit{inter-} + \textit{nation}）、\textbf{intervene}（干预；\textit{inter-} + \textit{venire} $\to$ 来到中间$\to$ 介入）、\textbf{interview}（面试；采访；\textit{inter-} + \textit{videre} $\to$ 互相看$\to$ 面谈）。
\end{origin}

\vw{outcome}{/ˈaʊtkʌm/}{n.}

\sense{n.}{结果，效果}
\col{the outcome of sth}
{The outcome of the election was a surprise to most observers.}
\col{outcome-based education 成果导向教育}
{Modern teaching methods focus more on outcomes than on the teaching process itself.}

\begin{gram}{outcome / result / consequence / aftermath 的"结果"细分}
\textbf{outcome}（结果——多指最终结局：\textit{the outcome of the match}）；\textbf{result}（结果——通用：\textit{the result of the experiment}）；\textbf{consequence}（后果——多指长远/重要：\textit{the consequences of one's actions}）；\textbf{aftermath}（余波——多指灾后的：\textit{the aftermath of the war}）。
\end{gram}

\begin{origin}
18 世纪由 \textit{out}（外）+ \textit{come}（来）合成。同源词 \textbf{outcome}、\textbf{income}（收入；\textit{in} 入 + come）、\textbf{overcome}（克服；\textit{over} 上 + come $\to$ 翻越过去）、\textbf{become}（成为；\textit{be-} 在旁 + come $\to$ 来到旁边$\to$变成）、\textbf{come}（来）、\textbf{comeback}（回归；反扑）。
\end{origin}

\vw{reef}{/riːf/}{n.}

\sense{n.}{礁，礁脉}
\col{a coral reef 珊瑚礁}
{The Great Barrier Reef is the world's largest coral reef system.}
\col{reef knot 平结}
{Tie the rope with a reef knot to make sure it holds.}

\begin{gram}{reef / shoal / bank / sandbar 的"浅水高地理"细分}
\textbf{reef}（礁——水下或近水面的岩石：\textit{coral reef}）；\textbf{shoal}（浅滩——沙子形成：\textit{a sand shoal}）；\textbf{bank}（浅滩——河/海中的：\textit{a sand bank}）；\textbf{sandbar}（沙洲——沙滩向海延伸的：\textit{a sandbar at the mouth of the river}）。reef 多为岩石，shoal/bank/sandbar 多为沙。
\end{gram}

\begin{origin}
16 世纪可能借自荷兰语 \textit{rif}（礁），与古诺斯语 \textit{rif}（肋；岩石）相关。同源词 \textbf{reef}、荷兰语 \textit{rijf}（骨架；肋骨）。reef 与 rib（肋）同源——礁石像海中浮出的"肋骨"。
\end{origin}

\vw{coral}{/ˈkɒrəl/}{n. \& adj.}

\sense{n.}{珊瑚}
\col{coral reef 珊瑚礁}
{The island is surrounded by some of the most diverse coral reefs in the world.}

\sense{adj.}{珊瑚色的}
\col{coral pink/lips}
{She wore a dress in a delicate coral pink.}

\begin{origin}
14 世纪借自古法语 \textit{coral}，拉丁语 \textit{corallum}，希腊语 \textit{korallion}（珊瑚）。词源不详——一说与希伯来语 \textit{goral}（石头）相关。同源词 \textbf{coral}、\textbf{coraline}（珊瑚的；珊瑚状的）。希腊人以为珊瑚是海底的"石化之花"。
\end{origin}

\vw{ecology}{/iˈkɒlədʒi/}{n.}

\sense{n.}{生态，生态学}
\col{the ecology of ...}
{The ecology of the rainforest is extremely complex.}
\col{ecological damage/diversity}
{The new road caused severe ecological damage to the surrounding area.}

\begin{origin}
19 世纪由希腊语 \textit{oikos}（家、住所）+ \textit{-logia}（学问）合成，1866 年由德国生物学家 Ernst Haeckel 创造。同源词 \textbf{ecology}、\textbf{ecological}（生态的）、\textbf{ecosystem}（生态系统）、\textbf{economy}（经济——希腊语 \textit{oikos} 家 + \textit{nomos} 法则$\to$ 治家之道）、\textbf{ecumenical}（全基督教的；\textit{oikos} 居住地$\to$ 共同的住所）。
\end{origin}

\vw{agenda}{/əˈdʒendə/}{n.}

\sense{n.}{计划，方案}
\col{an agenda for sth}
{We need a clear agenda for the next five years.}
\col{a hidden agenda 隐藏的议程}
{I suspected the director had a hidden agenda when he proposed the merger.}

\sense{n.}{议程表，议事日程}
\col{on the agenda}
{The question of climate change is high on the agenda of next week's summit.}
\col{item on the agenda}
{Let's move to the next item on the agenda.}

\begin{gram}{agenda / schedule / itinerary / programme 的"程"细分}
\textbf{agenda}（议程——多指会议：\textit{on the agenda}）；\textbf{schedule}（时间表——多指工作/交通：\textit{a train schedule}）；\textbf{itinerary}（行程——多指旅行：\textit{a travel itinerary}）；\textbf{programme}（节目单；计划——通用：\textit{a concert programme}）。
\end{gram}

\begin{origin}
17 世纪借自拉丁语 \textit{agenda}（待办之事），是 \textit{agere}（做）的中性复数。同源词 \textbf{agent}（代理人；动作者）、\textbf{agency}（机构）、\textbf{act}（行动）、\textbf{action}（行动）、\textbf{active}（积极的）、\textbf{react}（反应）、\textbf{transact}（办理；交易）、\textbf{exact}（精确的；\textit{ex-} 出 + agere $\to$ 强求出来）。agenda 字面"应该做的事"——议程是"需要被做之事"。
\end{origin}

\vw{pagoda}{/pəˈɡəʊdə/}{n.}

\sense{n.}{佛塔}
\col{a Buddhist pagoda}
{The ancient Buddhist pagoda has stood on the hill for over a thousand years.}

\begin{origin}
16 世纪借自葡萄牙语 \textit{pagode}，可能源自梵语 \textit{bhagavati}（神圣的；幸运的；Bhagavati 神），最终追溯至梵语 \textit{bhaga}（幸运的；分配者）。"pagoda" 的欧洲传播路线：印度 → 葡萄牙语 → 英语。中文"塔"是外来词——与梵语 \textit{sthupa} 有关。
\end{origin}

\vw{nail}{/neɪl/}{n. \& vt.}

\sense{n.}{钉子}
\col{drive/hit a nail 钉钉子}
{She drove the nail into the wood with three firm hits of the hammer.}
\col{a nail in the coffin 棺材上的钉子（致命打击）}
{The latest scandal was another nail in the coffin of his political career.}

\sense{n.}{指甲}
\col{finger/toe nail}
{Keep your finger nails short and clean.}

\sense{vt.}{固定}
\col{nail sth down 钉牢；确定}
{Let's nail down the details of the contract before we sign it.}
\col{nail sth up 把……钉起来}
{The broken windows were nailed up with boards.}

\begin{gram}{nail / screw / bolt / pin 的"钉固件"细分}
\textbf{nail}（钉子——木制：\textit{drive a nail}）；\textbf{screw}（螺丝——金属螺纹：\textit{turn a screw}）；\textbf{bolt}（螺栓——大号：\textit{an anchor bolt}）；\textbf{pin}（别针；销：\textit{a safety pin}）。nailed down 比 "fixed" 更口语化——"确定"如"钉牢"。
\end{gram}

\begin{origin}
古英语 \textit{nægl}（指甲；钉子），原始日耳曼语 *\textit{naglaz}。同源词 \textbf{nail}、\textbf{nail-biting}（咬指甲——焦虑的）、德语 \textit{Nagel}（钉子）、荷兰语 \textit{nagel}（钉子；指甲）。nagel 与英语 \textit{nag}（唠叨）形近但实不同源。
\end{origin}

\vw{panel}{/ˈpænl/}{n.}

\sense{n.}{专家咨询组，讨论小组}
\col{a panel of experts}
{A panel of experts will discuss the issue on tonight's programme.}
\col{panel discussion 专题讨论}
{She took part in a panel discussion on education reform.}

\sense{n.}{嵌板}
\col{a wooden/glass panel}
{The room has beautiful oak panels on the walls.}

\sense{n.}{仪表盘}
\col{control panel 控制面板}
{The control panel is located at the front of the aircraft.}

\begin{origin}
14 世纪中古英语 \textit{panel}（一块布；长方形布片），古法语 \textit{panel}（小方块布；布片），源于拉丁语 \textit{pannus}（布；衣料）。同源词 \textbf{panel}、\textbf{panelling}（镶板）、\textbf{empanel}（列入陪审团名单；\textit{em-} 入 + panel $\to$ 进入名单）、\textbf{panelist}（讨论小组成员）、\textbf{pane}（窗格玻璃——\textit{pannus} 布的延伸）。"panel" 原作"一块布"，现代义"专家小组"来自 16 世纪陪审员"列成一排坐"——像挂墙上的"镶板"。
\end{origin}

\vw{specialist}{/ˈspeʃəlɪst/}{n. \& adj.}

\sense{n.}{专家}
\col{a specialist in/on sth}
{She is a specialist in Renaissance art.}
\col{consult a specialist}
{If the symptoms persist, consult a specialist.}

\sense{n.}{专科医生}
\col{a heart/cancer/eye specialist}
{After the accident, she was referred to a heart specialist.}

\sense{adj.}{专业的}
\col{specialist advice/knowledge}
{For specialist advice on tax matters, consult a professional accountant.}

\begin{gram}{specialist / expert / authority / professional 的"专家"细分}
\textbf{specialist}（专家——专攻某一领域：\textit{a heart specialist}）；\textbf{expert}（能手；行家——多指经验/技能：\textit{a computer expert}）；\textbf{authority}（权威——多用于学术/官方：\textit{an authority on Chinese history}）；\textbf{professional}（专业人士——泛指）。specialist 偏医学/学科，expert 偏技能。
\end{gram}

\begin{origin}
19 世纪由 \textit{special}（特别的）+ \textit{-ist}（……的人）合成。special 14 世纪借自古法语 \textit{especial}，拉丁语 \textit{specialis}（特别的；专门的），源于 \textit{species}（种类；外观）。同源词 \textbf{special}、\textbf{specialist}、\textbf{specialize}（专攻）、\textbf{specialty}（专长；特产）、\textbf{species}（物种）、\textbf{specify}（指定；\textit{species} + \textit{-fy} $\to$ 使成为特定种类）、\textbf{especially}（尤其）、\textbf{especial}（特别的——古雅）。
\end{origin}

\vw{assess}{/əˈses/}{vt.}

\sense{vt.}{评估，评定；估算}
\col{assess sth}
{Teachers assess students' progress through a combination of tests and projects.}
\col{assess the damage/situation}
{It is too early to assess the full damage caused by the earthquake.}

\begin{gram}{assess / estimate / evaluate / appraise 的"估"细分}
见 evaluate 条。\textbf{assess}（评估——多用于税务/考试/局势：\textit{assess the value}）；\textbf{estimate}（粗略估计——多用于数量/价格：\textit{estimate the cost}）；\textbf{evaluate}（综合判断——多用于方案/表现：\textit{evaluate the proposal}）；\textbf{appraise}（专业评估——多用于资产/房产：\textit{appraise the property}）。
\end{gram}

\begin{origin}
15 世纪借自古法语 \textit{assesser}，拉丁语 \textit{assessare}（对……征税；评估税额），源于 \textit{assidere}（坐在旁边；辅助）：\textit{ad-}（朝向）+ \textit{sedere}（坐）。同源词 \textbf{assess}、\textbf{assessment}（评估）、\textbf{assessor}（评估人；税务稽查员）、\textbf{assiduous}（勤勉的；\textit{assidere} 坐在旁边辅助$\to$ 勤勉的）、\textbf{sedan}（轿车——\textit{sedere} 坐$\to$ 乘坐之物）、\textbf{sedentary}（久坐的；\textit{sedere} 坐$\to$ 久坐的生活）、\textbf{reside}（居住；\textit{re-} 回 + sedere $\to$ 坐定$\to$ 居住）、\textbf{preside}（主持；\textit{prae-} 前 + sedere $\to$ 坐在前面$\to$ 主持）、\textbf{president}（总统；主席——\textit{prae-} 前 + sedere $\to$ 坐在前面的人）。
\end{origin}

\vw{collapse}{/kəˈlæps/}{vi. \& n.}

\sense{vi.}{倒塌，坍塌}
\col{collapse under the weight/pressure}
{The old bridge collapsed under the weight of the heavy trucks.}
\col{collapse in a heap}
{She collapsed in a heap at the foot of the stairs.}

\sense{vi.}{昏倒}
\col{collapse from exhaustion}
{He collapsed from exhaustion after running the marathon.}

\sense{n.}{昏倒}
\col{a state of collapse}
{The patient was rushed to hospital in a state of collapse.}

\sense{n.}{倒塌}
\col{the collapse of a building/ceiling}
{Investigators are still trying to determine the cause of the building's collapse.}

\begin{gram}{collapse / fall down / topple / cave in 的"倒"细分}
\textbf{collapse}（塌——多指结构性失效：\textit{the roof collapsed}）；\textbf{fall down}（倒下——通用：\textit{fall down the stairs}）；\textbf{topple}（倾倒——多指不稳而翻倒：\textit{the statue toppled over}）；\textbf{cave in}（内陷——多指地面/天花板：\textit{the floor caved in}）。
\end{gram}

\begin{origin}
18 世纪借自拉丁语 \textit{collapsus}（倒塌、衰落），是 \textit{collabi}（一起倒下、倒塌）的过去分词：\textit{col-}（com- 同化形式"一起"）+ \textit{labi}（滑、倒）。同源词 \textbf{collapse}、\textbf{collapsible}（可折叠的）、\textbf{lapse}（失误；流逝；\textit{labi} $\to$ 滑走）、\textbf{elapse}（消逝；\textit{e-} 出 + labi $\to$ 滑出$\to$ 时间流逝）、\textbf{collapse of a deal}（谈判破裂）、\textbf{collapse of a company}（公司倒闭）。
\end{origin}

\vw{consultation}{/ˌkɒnslˈteɪʃn/}{n.}

\sense{n.}{咨询，商讨}
\col{in consultation with sb}
{The decision was taken in consultation with senior staff members.}
\col{undergo consultation}
{The new policy underwent extensive consultation with stakeholders.}

\sense{n.}{就诊}
\col{a medical consultation}
{If symptoms persist, book a medical consultation with your doctor.}

\begin{origin}
15 世纪借自拉丁语 \textit{consultationem}（商讨、咨询），源于 \textit{consultare}（商议、咨询），是 \textit{consulere}（商量、咨询）的反复形式。同源词 \textbf{consult}（咨询；查阅）、\textbf{consultant}（顾问；会诊医生）、\textbf{consultation}（咨询）、\textbf{consultative}（咨询的；协商的）、\textbf{counsel}（建议；商议；\textit{consulere} 缩减）、\textbf{council}（议会；\textit{consilium} 会议）、\textbf{reconcile}（调和；\textit{re-} 再 + \textit{con-} 一起 + \textit{calare} 召集$\to$ 再召集$\to$ 调和）。
\end{origin}

\vw{suspend}{/səˈspend/}{vt.}

\sense{vt.}{暂停}
\col{suspend sth}
{The game was suspended because of the heavy rain.}
\col{suspend sb from sth}
{He was suspended from school for two weeks.}

\sense{vt.}{悬，挂}
\col{suspend sth from/on sth}
{A large lamp was suspended from the ceiling.}
\col{suspended in sth}
{The dust was suspended in the air for hours after the explosion.}

\begin{gram}{suspend / interrupt / halt / pause 的"停"细分}
\textbf{suspend}（暂停——可恢复：\textit{suspend a meeting} 暂时休会）；\textbf{interrupt}（打断——常指他人：\textit{interrupt a conversation}）；\textbf{halt}（停止——多指完全停止：\textit{halt production} 停产）；\textbf{pause}（暂停——主动的：\textit{pause for breath} 喘口气）。
\end{gram}

\begin{origin}
15 世纪借自拉丁语 \textit{suspendere}（悬挂；暂停）：\textit{sus-}（sub- 同化形式"在下"）+ \textit{pendere}（悬挂、称重）。同源词 \textbf{suspend}、\textbf{suspender}（吊带；\textit{suspenders} 袜带）、\textbf{suspense}（悬念；\textit{suspendere} 悬挂$\to$ 悬而未决$\to$ 悬念）、\textbf{suspension}（暂停；悬挂）、\textbf{append}（附加；\textit{ap-} 朝 + \textit{pendere} $\to$ 挂在上面$\to$ 附加）、\textbf{pendulum}（钟摆）、\textbf{perpendicular}（垂直的；\textit{per-} 贯穿 + \textit{pendere} $\to$ 垂直悬挂的）、\textbf{pensive}（沉思的；\textit{pendere} 悬挂$\to$ 悬挂不定$\to$ 沉思的）。
\end{origin}

\vw{division}{/dɪˈvɪʒn/}{n.}

\sense{n.}{分歧}
\col{a division between/within ...}
{There is a deep division within the party over the issue of European integration.}
\col{cause/create division}
{The war caused serious divisions within the country.}

\sense{n.}{分开，分隔}
\col{the division of sth into sth}
{The division of the country into two states was a long and painful process.}
\col{division of labour 劳动分工}
{Modern manufacturing depends on the division of labour among specialised workers.}

\begin{origin}
14 世纪借自拉丁语 \textit{divisionem}（分开、分配），源于 \textit{dividere}（分开、分割）：\textit{di-}（dis- 同化形式"分开"）+ \textit{videre}（分——与 \textit{videre} 看不同源；实为 \textit{weid-} 分）。同源词 \textbf{divide}（分；除）、\textbf{division}（部门；分歧）、\textbf{divisor}（除数）、\textbf{individual}（个体的；\textit{in-} 不 + \textit{dividere} $\to$ 不可再分的）、\textbf{subdivide}（细分）、\textbf{widow}（寡妇——\textit{weid-} 分$\to$ 被分离的女人）。
\end{origin}

\vw{architect}{/ˈɑːkɪtekt/}{n.}

\sense{n.}{建筑师}
\col{an architect of/for sth}
{The architect designed several famous buildings in the city centre.}

\sense{n.}{设计师}
\col{the architect of a plan/policy}
{She was the chief architect of the new education policy.}

\begin{origin}
16 世纪借自拉丁语 \textit{architectus}，希腊语 \textit{arkhitektōn}（建筑大师；总工匠）：\textit{arkhi-}（首、主要）+ \textit{tekton}（工匠；建造者）。同源词 \textbf{architect}、\textbf{architecture}（建筑学；建筑风格）、\textbf{architectural}（建筑的）、\textbf{arch-}（首、主要——\textit{archbishop} 大主教、\textit{archenemy} 主要敌人）。"architect" 字面"主要工匠"——古希腊人对建造者的高级称谓。
\end{origin}

\vw{operator}{/ˈɒpəreɪtə(r)/}{n.}

\sense{n.}{操作人员，技工}
\col{a machine operator 机器操作员}
{She works as a machine operator in a factory.}

\sense{n.}{接线员}
\col{a telephone operator 电话接线员}
{At night, the only people you can call are the emergency operator.}

\begin{gram}{operator / worker / labourer 的"工"细分}
\textbf{operator}（操作员——多指机器/系统：\textit{a computer operator}）；\textbf{worker}（工人——通用：\textit{a factory worker}）；\textbf{labourer}（体力劳动者——多指户外：\textit{a farm labourer}）；\textbf{craftsman}（工匠——多指手工艺：\textit{a skilled craftsman}）。operator 强调"操控机器/系统"。
\end{gram}

\begin{origin}
17 世纪由 \textit{operate}（操作）+ \textit{-or} 派生。operate 17 世纪借自拉丁语 \textit{operari}（工作、操作），源于 \textit{opus}（工作；作品）。同源词 \textbf{operate}、\textbf{operation}（操作；手术）、\textbf{operator}、\textbf{operate}、\textbf{cooperate}（合作；\textit{co-} 一起 + operari $\to$ 一起工作）、\textbf{co-operation}、\textbf{opus}（作品；\textit{magnum opus} 杰作）、\textbf{opera}（歌剧——意大利语 \textit{opera} $\to$ 拉丁语 \textit{opera} 工作的复数$\to$ 音乐作品）。
\end{origin}

\vw{incident}{/ˈɪnsɪdənt/}{n.}

\sense{n.}{事件}
\col{an incident at/in ...}
{The police are investigating an incident at the embassy last night.}

\sense{n.}{冲突}
\col{a diplomatic/border incident}
{The two countries were on the brink of war after a serious border incident.}

\begin{gram}{incident / event / occurrence / happening 的"事件"细分}
\textbf{incident}（事件——多指不寻常/可叙述的：\textit{a strange incident}）；\textbf{event}（事件——多指重要/有意义的：\textit{historical events}）；\textbf{occurrence}（发生；偶发事件——偏中性）；\textbf{happening}（发生；偶发事件——口语化；20 世纪 60 年代有"Happening"艺术流派）。incident 偏戏剧性。
\end{gram}

\begin{origin}
15 世纪借自拉丁语 \textit{incidentem}（偶然发生的；降临的），是 \textit{incidere}（降临、发生、落下）的现在分词：\textit{in-}（在上）+ \textit{cadere}（落）。同源词 \textbf{incident}、\textbf{incidental}（附带的；偶然的）、\textbf{incidentally}（顺便说一下）、\textbf{coincide}（同时发生；\textit{co-} 一起 + \textit{incidere} $\to$ 同时落下$\to$ 重合）、\textbf{coincidence}（巧合）、\textbf{accident}（事故——\textit{ac-} 朝 + \textit{cadere} $\to$ 偶然降临$\to$ 事故）、\textbf{decadent}（颓废的——\textit{de-} 下 + \textit{cadere} $\to$ 落下的$\to$ 衰落的）。
\end{origin}

\vw{cooperation}{/kəʊˌpɒpəˈreɪʃn/}{n.}

\sense{n.}{合作，协作}
\col{in cooperation with sb}
{The two universities are working in close cooperation on the project.}
\col{international cooperation}
{International cooperation is essential to solving global problems like climate change.}

\sense{n.}{协助，配合}
\col{with the cooperation of sb}
{The investigation was carried out with the full cooperation of the local police.}

\begin{origin}
15 世纪由 \textit{cooperate}（合作）+ \textit{-ion} 派生。cooperate 17 世纪借自拉丁语 \textit{cooperari}（一起工作）：\textit{co-}（一起）+ \textit{operari}（工作）。同源词 \textbf{cooperate}、\textbf{cooperation}、\textbf{cooperative}（合作的；合作社）、\textbf{operate}、\textbf{operation}、\textbf{operator}、\textbf{opus}（作品）、\textbf{opera}（歌剧——意大利语 \textit{opera} 拉丁语 \textit{opera} 工作的复数）。"cooperate" 字面"一起工作"——合作即共事。
\end{origin}

\vw{forecast}{/ˈfɔːkɑːst/}{vt. \& n.}

\sense{vt.}{预测，预报}
\col{forecast sth}
{Economists have forecast a sharp rise in inflation next year.}
\col{forecast that ...}
{The weather service forecast that heavy rain would continue throughout the week.}

\sense{n.}{预测，预报}
\col{a weather/sales forecast}
{According to the weather forecast, snow is expected overnight.}
\col{a forecast for sth}
{The company's annual sales forecast has been revised upwards.}

\begin{gram}{forecast 的过去式与过去分词：forecast 或 forecasted}
\textbf{forecast} 的过去式与过去分词既可保持原形 \textit{forecast}（不规则），也可加 \textit{-ed} 变 \textit{forecasted}。英式英语偏好不规则；美式英语多见 \textit{forecasted}。两者皆正确。
\end{gram}

\begin{origin}
15 世纪由 \textit{fore-}（预先）+ \textit{cast}（投、抛；计算）合成。cast 在此是古义"计算"——占卜/预测。同源词 \textbf{forecast}、\textbf{forefather}（祖先；\textit{fare} 前 + \textit{father}）、\textbf{foresee}（预见；\textit{fare} 前 + \textit{see} 看）、\textbf{forecastle}（船首楼；\textit{fare} 前 + \textit{castle} 城堡$\to$ 前方城堡——船长指挥塔）、\textbf{before}（在……之前）。
\end{origin}

\vw{interrupt}{/ˌɪntəˈrʌpt/}{vt. \& vi.}

\sense{vt.}{打扰，插嘴}
\col{interrupt sb}
{Please don't interrupt me while I'm speaking.}
\col{interrupt sb with sth}
{She kept interrupting him with silly questions.}

\sense{vt.}{使暂停，使中断}
\col{interrupt sth}
{The snowstorm interrupted all flights at the airport.}
\col{interrupt the flow/progress of sth}
{The war interrupted the smooth flow of trade between the two countries.}

\sense{vi.}{打扰，插嘴}
\col{interrupt (with sth)}
{It's rude to interrupt when someone else is talking.}

\begin{origin}
15 世纪借自拉丁语 \textit{interruptus}（打断、中断），是 \textit{interrumpere}（打断、破裂）的过去分词：\textit{inter-}（在……之间）+ \textit{rompere}（打破、破裂）。同源词 \textbf{interrupt}、\textbf{interruption}（中断）、\textbf{interrupted}（被打断的）、\textbf{rupture}（破裂）、\textbf{disrupt}（扰乱；\textit{dis-} 离 + \textit{rompere}）、\textbf{erupt}（爆发）、\textbf{abrupt}（突然的；\textit{ab-} 离 + \textit{rompere} $\to$ 突然断裂）。
\end{origin}

\vw{conservation}{/ˌkɒnsəˈveɪʃn/}{n.}

\sense{n.}{保护}
\col{wildlife/nature conservation}
{She works for a wildlife conservation organisation.}
\col{energy conservation 能源保护}
{Energy conservation is becoming increasingly important as resources dwindle.}

\sense{n.}{保持}
\col{conservation of momentum/energy 动量守恒/能量守恒}
{The law of conservation of energy is one of the fundamental principles of physics.}

\begin{origin}
15 世纪由 \textit{conserve}（保存）+ \textit{-ation} 派生。conserve 14 世纪借自古法语 \textit{conserver}，拉丁语 \textit{conservare}（保存、保护）：\textit{con-}（完全）+ \textit{servare}（保存；保护）。同源词 \textbf{conserve}（节约；保护）、\textbf{conservation}、\textbf{conservative}（保守的；保守党）、\textbf{preserve}（保存）、\textbf{reserve}（保留；预订）、\textbf{observe}（观察；遵守；\textit{ob-} 朝 + servare $\to$ 面向保护$\to$ 观察/遵守）、\textbf{deserve}（应得）。
\end{origin}

\vw{motion}{/ˈməʊʃn/}{n.}

\sense{n.}{提议，动议}
\col{a motion to do sth}
{She proposed a motion to postpone the meeting until next week.}
\col{carry/pass a motion 通过动议}
{The motion was carried by a large majority.}

\sense{n.}{运动}
\col{in motion}
{The train was already in motion when she jumped aboard.}
\col{laws of motion 运动定律}
{Newton's laws of motion form the foundation of classical mechanics.}

\sense{n.}{动作}
\col{graceful hand motions}
{He described the whole scene with broad, sweeping hand motions.}

\begin{gram}{motion / movement / move / mobility 的"动"细分}
\textbf{motion}（运动——物理/机械意义：\textit{in motion} 运动中）；\textbf{movement}（运动；动向——多指社会/政治：\textit{a political movement}）；\textbf{move}（动作；步骤——通用：\textit{one's next move}）；\textbf{mobility}（流动性；可动性——多形容阶层/位置：\textit{social mobility}）。
\end{gram}

\begin{origin}
15 世纪借自拉丁语 \textit{motionem}（运动；动议），源于 \textit{motus}（移动），是 \textit{movere}（移动、推动）的过去分词。同源词 \textbf{motion}、\textbf{motivate}（激发）、\textbf{motive}（动机）、\textbf{motor}（发动机）、\textbf{emotion}（情绪；\textit{e-} 出 + movere $\to$ 移出心外）、\textbf{promote}（促进；\textit{pro-} 前 + movere $\to$ 向前推动）、\textbf{remote}（遥远的；\textit{re-} 离 + movere $\to$ 移开$\to$ 远的）。
\end{origin}

\vw{vote}{/vəʊt/}{vt. \& vi. \& n.}

\sense{vt. \& vi.}{投票，表决}
\col{vote for/against sb/sth}
{A majority of members voted for the new proposal.}
\col{vote to do sth}
{The committee voted to accept the recommendation.}

\sense{n.}{选票；投票}
\col{cast a vote 投票}
{Over 60\% of the eligible voters cast their vote in the election.}
\col{a vote of thanks 鸣谢}
{The chairman proposed a vote of thanks to the organisers.}

\begin{gram}{vote / ballot / poll / referendum 的"投票"细分}
\textbf{vote}（投票——泛指：\textit{cast a vote}）；\textbf{ballot}（无记名投票——多指选票本身：\textit{by secret ballot}）；\textbf{poll}（民意调查；投票处：\textit{an opinion poll}）；\textbf{referendum}（公投——就某一具体议题）。vote 偏动作，ballot 偏形式，poll 偏调查/统计，referendum 偏全民议题。
\end{gram}

\begin{origin}
15 世纪借自拉丁语 \textit{votum}（誓言、愿望、祈祷），源于 \textit{vovere}（许愿、宣誓）。同源词 \textbf{vote}、\textbf{voter}（选民）、\textbf{votive}（奉献的；还愿的；\textit{a votive candle} 还愿蜡烛）、\textbf{vow}（誓言；\textit{vovere} 的英语化）、\textbf{devote}（奉献；\textit{de-} 加强 + \textit{vovere} $\to$ 郑重许愿$\to$ 奉献）、\textbf{devout}（虔诚的；\textit{de-} 加强 + \textit{votum} $\to$ 极度虔诚的）。"vote" 字面"许下的愿"——古罗马人投票时向神许下庄严誓言。
\end{origin}

\vw{migration}{/maɪˈɡreɪʃn/}{n.}

\sense{n.}{迁徙，迁移，移居}
\col{bird migration 鸟类迁徙}
{The wetlands are an important stopping point for birds on their annual migration.}
\col{mass migration 大规模迁徙}
{Mass migration from rural areas to cities has been a feature of modern life.}
\col{net migration 净迁移}
{The region has experienced positive net migration over the past decade.}

\begin{gram}{migration / immigration / emigration 的"迁"细分}
\textbf{migration}（迁移——泛指：\textit{human migration}）；\textbf{immigration}（移入——从国外进入：\textit{immigration policy}）；\textbf{emigration}（移出——离开本国：\textit{emigration from Europe}）。三者的逻辑：migration 涵盖 imm/em 两者，im- 进 vs. e- 出。
\end{gram}

\begin{origin}
17 世纪由 \textit{migrate}（迁徙）+ \textit{-ion} 派生。migrate 17 世纪借自拉丁语 \textit{migrare}（迁移、移居），可能与希腊语 \textit{ameibein}（交换）相关。同源词 \textbf{migrate}、\textbf{migration}、\textbf{migrant}（迁徙者；移民）、\textbf{emigrate}（移出）、\textbf{immigrate}（移入）、\textbf{transmigrate}（转世；\textit{trans-} 跨越 + migrare）。"migrate" 与 "change" 在古印欧语中同源——迁移即是改变。
\end{origin}

\vw{legacy}{/ˈleɡəsi/}{n.}

\sense{n.}{遗产}
\col{a legacy of sth}
{The war left a legacy of destruction and bitterness that lasted for generations.}
\col{leave a legacy}
{She left a legacy of kindness and generosity that her children still cherish.}
\col{legacy system 遗留系统}
{The new software will eventually replace the company's ageing legacy systems.}

\begin{origin}
14 世纪中古英语 \textit{legacie}，古法语 \textit{legacie}（大使职责；遗产），中世纪拉丁语 \textit{legatia}（大使职责），源于 \textit{legatus}（被任命的；大使）。同源词 \textbf{legacy}、\textbf{legate}（教皇使节；特使）、\textbf{legation}（公使馆）、\textbf{relegate}（降级；\textit{re-} 回 + \textit{legare} $\to$ 遣回原位$\to$ 降级）、\textbf{legal}（法律的；与 legacy 形近但实不同源——legal 源自 \textit{lex} 法）。"legacy" 字面"使节留下之物"——特使死后留下的财产。
\end{origin}

\vw{splendour}{/ˈsplendə(r)/}{n.}

\sense{n.}{壮丽，雄伟}
\col{the splendour of ...}
{The splendour of the ancient temples left us speechless.}
\col{in all its splendour}
{The palace has been restored to all its former splendour.}

\begin{lex}{splendour / magnificence / glory / grandeur 的"壮"细分}
\textbf{splendour}（壮丽——偏视觉/壮美：\textit{the splendour of the sunset}）；\textbf{magnificence}（宏伟——偏规模/华丽：\textit{the magnificence of the palace}）；\textbf{glory}（光荣——偏成就/名声：\textit{the glory of the empire}）；\textbf{grandeur}（宏伟——偏气势/庄严：\textit{the grandeur of the Alps}）。
\end{lex}

\begin{origin}
15 世纪借自古法语 \textit{splendour}，拉丁语 \textit{splendor}（光辉、壮丽），源于 \textit{splendere}（发光、闪耀）。同源词 \textbf{splendid}（壮丽的；极好的）、\textbf{splendour}、\textbf{resplendent}（光辉灿烂的；\textit{re-} 加强 + splendere）。同源字根 \textit{splend-}（发光）——壮丽源自"发光之物"。
\end{origin}

\vw{reserve}{/rɪˈzɜːv/}{n. \& vt.}

\sense{n.}{自然保护区}
\col{a nature/wildlife/game reserve}
{The Serengeti National Park is one of the world's most famous wildlife reserves.}

\sense{n.}{储备}
\col{reserves of ...}
{The country has large reserves of oil and natural gas.}
\col{keep/have sth in reserve}
{We always keep some money in reserve for emergencies.}

\sense{vt.}{预订；预留}
\col{reserve sth for sb/sth}
{I'd like to reserve a table for two at the restaurant for 8 o'clock.}

\sense{vt.}{保留}
\col{reserve the right to do sth}
{The company reserves the right to refuse service to anyone.}

\begin{gram}{reserve / reservation 的可数性}
\textbf{reserve}（n.）作"储备"时常用复数 \textit{reserves}（\textit{oil reserves}）；作"自然保护区"时用单数（\textit{a nature reserve}）；作"预订"时用单数（\textit{make a reserve} 较少用）。\textbf{reservation}（n. 预订；保留意见）作"预订"时常单数（\textit{make a reservation}）；作"保留意见"时常复数（\textit{have reservations about...}）。注意英美差异：英式 \textit{booking} = 美式 \textit{reservation}。
\end{gram}

\begin{origin}
14 世纪借自古法语 \textit{reserver}，拉丁语 \textit{reservare}（保留、储存）：\textit{re-}（回）+ \textit{servare}（保存、保护）。同源词 \textbf{reserve}、\textbf{reservation}、\textbf{reservoir}（水库；\textit{reservoir} 法语形式$\to$ 储存之处）、\textbf{preserve}（保存）、\textbf{conserve}（保护）、\textbf{observe}（观察；遵守）、\textbf{deserve}（应得）。
\end{origin}

\vw{basin}{/ˈbeɪsn/}{n.}

\sense{n.}{盆地}
\col{the Amazon Basin}
{The Amazon Basin covers a vast area of South America.}

\sense{n.}{洗脸盆}
\col{a wash basin}
{She filled the wash basin with hot water.}

\sense{n.}{盆}
\col{a basin of water}
{He mixed the flour with a basin of cold water.}

\sense{n.}{流域}
\col{the Thames Basin}
{All the rivers in the Thames Basin flow into the Thames.}

\begin{origin}
13 世纪借自古法语 \textit{bacin}（盆、水槽），可能与中世纪拉丁语 \textit{bacchinus}（水盆）相关。同源词 \textbf{basin}、\textbf{basinful}（一盆之量）、德语 \textit{Becken}（盆；骨盆）。basin 与 base（底部）形近但实不同源。
\end{origin}

\vw{intense}{/ɪnˈtens/}{adj.}

\sense{adj.}{很大的，十分强烈的}
\col{intense heat/cold/pain}
{The intense heat of the desert made it impossible to travel by day.}
\col{intense competition/pressure}
{The competition for places at the university is extremely intense.}

\sense{adj.}{有强烈感情的，热切的}
\col{intense emotion/feeling/love}
{She felt an intense longing to return to her childhood home.}

\sense{adj.}{激烈的}
\col{an intense debate/gaze}
{He fixed her with an intense gaze that made her uncomfortable.}

\begin{gram}{intense / intensive / intent 的派生分叉}
\textbf{intense}（强烈的——形容程度：\textit{intense heat}）；\textbf{intensive}（密集的；彻底的——形容方法/过程：\textit{an intensive course} 强化课程；\textit{intensive care} 重症监护）；\textbf{intent}（专注的；坚决的——\textit{intent on doing}）。intense 偏物理/情感强度，intensive 偏时间/精力的密集投入。
\end{gram}

\begin{origin}
15 世纪借自中古法语 \textit{intens}，拉丁语 \textit{intensus}（绷紧的、强烈的），是 \textit{intendere}（伸向、指向、打算）的过去分词：\textit{in-}（向内）+ \textit{tendere}（伸、拉伸）。同源词 \textbf{intense}、\textbf{intensity}（强度）、\textbf{intensify}（加强）、\textbf{intensive}（密集的）、\textbf{intent}（意图；专注的）、\textbf{pretend}（假装；\textit{prae-} 在前 + tendere $\to$ 把意图放在前面）。
\end{origin}

\vw{primitive}{/ˈprɪmətɪv/}{adj.}

\sense{adj.}{简陋的，落后的}
\col{primitive weapons/tools}
{The explorers survived with only primitive tools and their own courage.}
\col{primitive conditions}
{The refugees were housed under primitive conditions in the mountains.}

\sense{adj.}{原始的，远古的}
\col{primitive man/society}
{Primitive man lived in caves and hunted for food.}
\col{primitive times/eras}
{The paintings depict scenes from primitive times.}

\begin{gram}{primitive / prime / primary / primal / primeval 的"原"分叉}
\textbf{primitive}（原始的；简陋的——多指早期人类/事物：\textit{primitive man}）；\textbf{prime}（首要的；最好的——\textit{prime minister} 首相；\textit{in one's prime} 壮年）；\textbf{primary}（最初的；主要的——\textit{primary school} 小学；\textit{primary concern}）；\textbf{primal}（最初的；本源的——\textit{primal scene} 原始场景）；\textbf{primeval}（远古的——\textit{primeval forest} 原始森林）。
\end{gram}

\begin{origin}
15 世纪借自拉丁语 \textit{primitivus}（最初的；最早的），源于 \textit{primus}（第一；最初），是 \textit{pro}（在前）的最高级。同源词 \textbf{primitive}、\textbf{prime}（最初的；最好的；\textit{primus} 英语化）、\textbf{primary}（最初的；主要的；\textit{primus} + \textit{-ary}）、\textbf{primate}（大主教；灵长类——\textit{primus} 第一$\to$ 最高位者）、\textbf{premier}（首相；最先的——古法语 \textit{primier}）、\textbf{prince}（王子——拉丁语 \textit{ princeps} 第一人）。"primitive" 字面"最初之物"——时间序列的最早端。
\end{origin}

\vw{humble}{/ˈhʌmbl/}{vt. \& adj.}

\sense{vt.}{使感到谦卑，使感到自惭}
\col{humble sb}
{The experience of working with the poor humbled him greatly.}

\sense{adj.}{谦虚的，谦卑的}
\col{be humble about sth}
{Despite her fame, she remains humble about her achievements.}
\col{in my humble opinion 依我拙见}
{In my humble opinion, the book is his best work to date.}

\begin{lex}{humble / modest / meek / lowly 的"谦"细分}
\textbf{humble}（谦卑的——含"知自己不足"：\textit{a humble hero}）；\textbf{modest}（谦虚的——多指不过分张扬：\textit{be modest about one's success}）；\textbf{meek}（温顺的——偏负面，暗示顺从过度：\textit{as meek as a lamb}）；\textbf{lowly}（卑微的——多形容出身/地位：\textit{lowly origins}）。humble 偏"自知之明"，modest 偏"不自我标榜"。
\end{lex}

\begin{origin}
13 世纪借自古法语 \textit{humble}，拉丁语 \textit{humilis}（贴近地面的；谦卑的），源于 \textit{humus}（地面；土壤）。同源词 \textbf{humble}、\textbf{humbly}（谦卑地）、\textbf{humility}（谦逊）、\textbf{humiliate}（羞辱；\textit{humilis} $\to$ 使贴近地面$\to$ 使受辱）、\textbf{humiliation}（屈辱）、\textbf{homo}（人——拉丁语 \textit{homo} 来自 \textit{humus} 土——"土地所生之物"$\to$ 人类）、\textbf{human}（人；人道的）、\textbf{humus}（腐殖质；\textit{humilis} 同一词根）。"humble" 字面"贴近地面"——谦卑是"将自己放在低处"。
\end{origin}

\vw{corridor}{/ˈkɒrɪdɔːr(r)/}{n.}

\sense{n.}{走廊，地带；通道}
\col{along/down the corridor}
{She walked along the corridor to the end room.}
\col{a corridor of power 权力走廊}
{Decisions are often made in the corridors of power, away from the public eye.}

\begin{origin}
16 世纪借自意大利语 \textit{corridore}（走廊；跑道），晚期拉丁语 \textit{currere}（跑）：\textit{cur-}（跑；行走）+ \textit{-idor} 后缀。同源词 \textbf{corridor}、\textbf{course}（跑；过程；\textit{currere} 的英语化）、\textbf{concourse}（汇合；广场；\textit{con-} 一起 + currere $\to$ 一起跑$\to$ 汇合）、\textbf{concur}（同时发生；同意）、\textbf{occur}（发生；\textit{oc-} 朝 + currere $\to$ 跑到面前$\to$ 出现）、\textbf{recur}（再发生；\textit{re-} 回 + currere）、\textbf{current}（当前的；水流；\textit{currere} 的现在分词）、\textbf{excursion}（远足；\textit{ex-} 出 + currere $\to$ 跑出去$\to$ 远行）。"corridor" 字面"跑道"——古罗马军营中士兵列队跑过的通道。
\end{origin}

\vw{pursue}{/pəˈsjuː/}{vt.}

\sense{vt.}{追求，致力于，执行}
\col{pursue a career/goal/dream}
{She decided to pursue a career in medicine after graduating.}
\col{pursue a policy/course of action}
{The government is pursuing a policy of economic reform.}

\sense{vt.}{追逐}
\col{pursue sb/sth}
{The police pursued the suspect through the city streets.}

\begin{lex}{pursue / chase / follow / track 的"追"细分}
\textbf{pursue}（追求；追捕——多指抽象/官方：\textit{pursue one's dream}）；\textbf{chase}（追逐——多指快速追赶：\textit{chase the thief}）；\textbf{follow}（跟随——通用：\textit{follow the path}）；\textbf{track}（追踪——多指循迹：\textit{track the animal}）。
\end{lex}

\begin{origin}
14 世纪借自古法语 \textit{pursuir}，拉丁语 \textit{prosequi}（追随、追踪）：\textit{pro-}（向前）+ \textit{sequi}（跟随）。同源词 \textbf{pursue}、\textbf{pursuit}（追求；追捕；\textit{in hot pursuit} 紧追不舍）、\textbf{prosecute}（起诉；\textit{pro-} 前 + \textit{sequi} $\to$ 追到前面$\to$起诉）、\textbf{consecutive}（连续的；\textit{con-} 一起 + \textit{sequi} $\to$ 一个接一个）、\textbf{sequence}（序列）、\textbf{consequence}（结果；\textit{con-} 一起 + \textit{sequi} $\to$ 随之而来的）。
\end{origin}

\vw{submit}{/səbˈmɪt/}{vt. \& vi.}

\sense{vt.}{提交，呈递}
\col{submit sth (to sb/sth)}
{All proposals must be submitted to the committee by Friday.}
\col{submit an application/report}
{She submitted her application for the job two weeks ago.}

\sense{vi.}{屈服}
\col{submit to sth}
{After a long siege, the defenders were forced to submit to the enemy.}
\col{refuse to submit to sb/sth}
{The workers refused to submit to the company's demands.}

\begin{gram}{submit / yield / surrender / succumb 的"屈服"细分}
\textbf{submit}（屈服——多指接受权威/规则：\textit{submit to the rules}）；\textbf{yield}（让步；屈服——多指不情愿的让步：\textit{yield to pressure}）；\textbf{surrender}（投降——多指军事/正式：\textit{surrender to the enemy}）；\textbf{succumb}（屈服于——多指疾病/诱惑：\textit{succumb to illness}）。submit 较温和，surrender/ succumb 较强。
\end{gram}

\begin{origin}
15 世纪借自拉丁语 \textit{submittere}（置于下方；屈服）：\textit{sub-}（在下）+ \textit{mittere}（送；放置）。同源词 \textbf{submit}、\textbf{submission}（提交；屈服）、\textbf{submissive}（顺从的）、\textbf{mission}（任务）、\textbf{missile}（导弹）、\textbf{emit}（发射）、\textbf{dismiss}（解雇）、\textbf{admit}（承认；\textit{ad-} 朝 + mittere $\to$ 允许入内）。"submit" 字面"放在下面"——屈服是把自我放在他人之下。
\end{origin}

\vw{emerge}{/ɪˈmɜːdʒ/}{vi.}

\sense{vi.}{出现}
\col{emerge from sth}
{The swimmer emerged from the water, gasping for breath.}
\col{emerge into the open 浮出水面；公开化}
{The truth finally emerged into the open after years of silence.}

\sense{vi.}{暴露}
\col{emerge (from sth)}
{New evidence has emerged from the investigation.}
\col{it emerges that ...}
{It emerged that the company had been losing money for years.}

\sense{vi.}{露头}
\col{emerge from ...}
{After the long winter, the first flowers began to emerge from the soil.}

\begin{origin}
17 世纪借自拉丁语 \textit{emergere}（浮出、出现）：\textit{e-}（出）+ \textit{mergere}（沉、浸）。同源词 \textbf{emerge}、\textbf{emergency}（紧急情况——\textit{emergere} 浮出$\to$ 突然出现$\to$ 紧急）、\textbf{emergent}（新兴的）、\textbf{submerge}（浸没；\textit{sub-} 在下 + mergere $\to$ 沉到下面$\to$ 浸没）、\textbf{immerse}（沉浸；\textit{im-} 入 + mergere $\to$ 浸入$\to$ 沉浸）、\textbf{merge}（合并）。"emerge" 字面"从水中浮出"——从淹没中现形。
\end{origin}

\vw{assign}{/əˈsaɪn/}{vt.}

\sense{vt.}{指定，指派}
\col{assign sb to sth}
{Two experienced officers were assigned to the case.}
\col{assign sb to do sth}
{I was assigned to interview the new applicants.}

\sense{vt.}{分配}
\col{assign sth to sb}
{The teacher assigned a different task to each student.}

\sense{vt.}{确定}
\col{assign a value/date}
{It is difficult to assign a value to the experience.}

\begin{gram}{assign / allocate / distribute / apportion 的"分配"细分}
\textbf{assign}（指定；指派——多指任务：\textit{assign homework}）；\textbf{allocate}（配给；划拨——多指资源/资金：\textit{allocate funds}）；\textbf{distribute}（分发；分配——多指物品：\textit{distribute leaflets}）；\textbf{apportion}（按比例分配——多指份额：\textit{apportion blame}）。assign 偏任务，allocate 偏资源，distribute 偏物品，apportion 偏份额。
\end{gram}

\begin{origin}
14 世纪借自古法语 \textit{assigner}，拉丁语 \textit{assignare}（标记、分配、转让）：\textit{ad-}（朝向、为）+ \textit{signare}（标记；签署），源于 \textit{signum}（标记；记号）。同源词 \textbf{assign}、\textbf{assignment}（任务）、\textbf{assignation}（约会；分配）、\textbf{sign}（标记；迹象）、\textbf{signal}（信号）、\textbf{design}（设计；\textit{de-} 下 + signare $\to$ 标出$\to$ 设计）、\textbf{resign}（辞职；\textit{re-} 回 + signare $\to$ 签字离开$\to$ 辞职）、\textbf{seal}（印章；\textit{signum} 缩小形式——封蜡印）。
\end{origin}

\vp{take shape}

\sense{短语}{成形，有了模样}
\ex{The building slowly began to take shape as the workers added more floors.}
\ex{After months of discussion, the plan is finally beginning to take shape.}

\begin{gram}{take shape / take form / take a turn / take a stand：take 短语"成形"族}
\textbf{take shape}（成形；\textit{the plan took shape}）；\textbf{take form}（变得具体；与 take shape 近义但更正式）；\textbf{take a turn}（发生变化：\textit{the situation took a turn for the worse} 形势恶化）；\textbf{take a stand}（表明立场：\textit{take a stand against sth} 反对某事）。
\end{gram}

\vw{string}{/strɪŋ/}{n.}

\sense{n.}{一批，一连串，一系列}
\col{a string of sth}
{A string of successes has marked his career.}
\col{a string of pearls/beads}
{She wore a string of pearls around her neck.}

\sense{n.}{线}
\col{a piece of string}
{Tie the package with a piece of string.}

\sense{n.}{弦}
\col{a violin/guitar string}
{The string broke as she was tuning the guitar.}

\begin{gram}{pull strings 拉关系}
\textbf{pull strings}（拉关系；走后门：\textit{He got the job by pulling strings}）、\textbf{no strings attached}（无附加条件：\textit{a loan with no strings attached}）、\textbf{string sb along}（吊着某人；愚弄）、\textbf{string sth out}（拖延）。string 抽象用法皆源自物理的"线/弦"——关系是"线"，条件是"绳"。
\end{gram}

\begin{origin}
古英语 \textit{streng}（线；绳子；弦），原始日耳曼语 *\textit{strangiz}。同源词 \textbf{string}、\textbf{stringy}（纤维质的；多筋的）、德语 \textit{Strang}（绳）、\textbf{stringent}（严格的；\textit{stringere} 拉紧——与 string 形近但实不同源，\textit{stringere} 拉紧是拉丁词）、\textbf{string bean}（菜豆——长如弦的豆）。
\end{origin}

\vw{cotton}{/ˈkɒtn/}{n.}

\sense{n.}{棉织物；棉，棉花}
\col{cotton wool 脱脂棉}
{She wrapped the wound in cotton wool.}
\col{cotton industry 棉纺业}
{The cotton industry was the backbone of the region's economy in the 19th century.}

\begin{lex}{cotton / wool / silk / linen 的"纤维"细分}
\textbf{cotton}（棉——植物纤维：\textit{a cotton shirt}）；\textbf{wool}（羊毛——动物纤维：\textit{a wool sweater}）；\textbf{silk}（丝绸——蚕丝：\textit{a silk dress}）；\textbf{linen}（亚麻——植物纤维：\textit{linen sheets}）。cotton 偏日常，wool 偏保暖，silk 偏奢华，linen 偏凉爽。
\end{lex}

\begin{origin}
14 世纪借自古法语 \textit{coton}，意大利语 \textit{cotone}，阿拉伯语 \textit{qutn}（棉花），可能最终源自埃及语。同源词 \textbf{cotton}、\textbf{cottonseed}（棉籽）、\textbf{cottony}（像棉的；柔软的）。cotton 经由阿拉伯贸易传入欧洲——丝绸之路的另一通道。
\end{origin}

\vw{Buddhist}{/ˈbʊdɪst/}{adj. \& n.}

\sense{adj.}{佛教的}
\col{a Buddhist monk/temple}
{We visited several ancient Buddhist temples during our trip to Southeast Asia.}

\sense{n.}{佛教徒}
\col{devout Buddhist}
{She has been a devout Buddhist for many years.}

\begin{origin}
17 世纪由 \textit{Buddha}（佛陀）+ \textit{-ist} 派生。Buddha 借自梵语 \textit{buddha}（觉悟者），源于动词词根 \textit{budh-}（觉悟、知道、醒）。同源词 \textbf{Buddha}、\textbf{Buddhism}（佛教）、\textbf{Buddhist}、\textbf{bodhi}（菩提——觉悟）。英语 \textbf{bode}（预示；\textit{bode well} 吉兆）可能与梵语 \textit{budh-} 共享印欧语根 *\textit{bheudh-}（觉醒）。
\end{origin}

\vw{expansion}{/ɪkˈspænʃn/}{n.}

\sense{n.}{扩展，扩大，膨胀}
\col{the expansion of ...}
{The expansion of the university over the past decade has been remarkable.}
\col{territorial expansion 领土扩张}
{The empire's territorial expansion reached its peak in the 19th century.}
\col{economic expansion 经济扩张}
{Economic expansion has created thousands of new jobs.}

\begin{gram}{expansion / expansionism / expansive 的派生分叉}
\textbf{expansion}（扩张——中性名词）；\textbf{expansionism}（扩张主义——含政治/军事色彩：\textit{colonial expansionism}）；\textbf{expansive}（广阔的；开朗的——形容词：\textit{an expansive view}；\textit{an expansive personality}）。注意 expansive 表"性格开朗"时与"扩张"无直接关联。
\end{gram}

\begin{origin}
17 世纪由 \textit{expand}（扩张）+ \textit{-sion} 派生。expand 17 世纪借自拉丁语 \textit{expandere}（展开、张开）：\textit{ex-}（出）+ \textit{pendere}（悬挂；称重）。同源词 \textbf{expand}、\textbf{expansion}、\textbf{expansive}、\textbf{expandable}（可扩展的）、\textbf{pace}（步；步伐；拉丁语 \textit{passus} 步，\textit{expandere} 的另一种词源引申——丈量土地时"伸开双臂"）、\textbf{passport}（护照——"通过的港口"）。"expand" 字面"张开"——从折叠/压缩状态打开。
\end{origin}

\vw{initiative}{/ɪˈnɪʃətɪv/}{n.}

\sense{n.}{倡议，新方案}
\col{a peace initiative 和平倡议}
{The government has launched a new peace initiative to end the long-running conflict.}
\col{on one's own initiative 主动地}
{She started the project on her own initiative, without waiting for approval.}

\sense{n.}{主动性}
\col{show initiative 表现出主动性}
{New employees are encouraged to show initiative in solving problems.}
\col{take the initiative 采取主动}
{Take the initiative in introducing yourself at the meeting.}

\begin{gram}{initiative / enterprise / drive 的"主动性"细分}
\textbf{initiative}（主动性；首创精神——多指个人：\textit{show initiative}）；\textbf{enterprise}（事业心；进取心——多指组织：\textit{private enterprise}）；\textbf{drive}（驱动力——多指内在推动：\textit{a person with drive}）。initiative 偏具体行动，enterprise 偏组织能力，drive 偏个人动力。
\end{gram}

\begin{origin}
17 世纪由 \textit{initiate}（开始；启动）+ \textit{-ive} 派生。initiate 16 世纪借自拉丁语 \textit{initiare}（开始、入门），源于 \textit{initium}（开始、入口），是 \textit{inire}（进入）的派生形式：\textit{in-}（入）+ \textit{ire}（走）。同源词 \textbf{initiate}、\textbf{initiation}（启动；入会仪式）、\textbf{initiative}、\textbf{initial}（最初的；\textit{initium} 最初）、\textbf{initially}（最初）、\textbf{itinerary}（行程；\textit{iter} 路——与 \textit{ire} 走同源）。
\end{origin}

\vw{showcase}{/ˈʃəʊkeɪs/}{n.}

\sense{n.}{展示的场合}
\col{a showcase for sth}
{The exhibition is a showcase for contemporary Chinese art.}
\col{showcase one's talents 展示才华}
{This is your chance to showcase your talents to a wider audience.}

\sense{n.}{玻璃柜台}
\ex{The shop's showcases displayed a glittering array of jewellery.}

\begin{origin}
20 世纪由 \textit{show}（展示）+ \textit{case}（盒子；柜）合成。同源词 \textbf{showcase}、\textbf{showcase} 的合成路径：show（展示；古英语 \textit{scēawian} 看）+ case（盒子；拉丁语 \textit{capere} 抓住$\to$ 容器）。合成词是现代英语构词的常用途径。
\end{origin}

\vw{elsewhere}{/ˌelsˈweə(r)/}{adv.}

\sense{adv.}{在（或去）别处}
\ex{If the product is not available in our shop, try elsewhere.}
\ex{She decided to look for a job elsewhere.}

\begin{gram}{elsewhere / somewhere / anywhere / nowhere 的复合副词}
四者由 \textit{else / some / any / no} + \textit{where} 合成。\textbf{elsewhere}（别处——特定之外：\textit{go elsewhere}）；\textbf{somewhere}（某处；某地：\textit{somewhere in the world}）；\textbf{anywhere}（任何地方——多用于疑问/否定：\textit{not anywhere}）；\textbf{nowhere}（无处——多用于否定：\textit{nowhere to be found}）。注意 elsewhere 已含 else 含义，不再说 \textit{else elsewhere}。
\end{gram}

\begin{origin}
古英语 \textit{elles}（其他；不同地）+ \textit{hwǣr}（哪里；何地）。同源词 \textbf{else}（其他的）、\textbf{elsewhere}、\textbf{where}（哪里）、\textbf{nowhere}、\textbf{somewhere}、\textbf{anywhere}。
\end{origin}

\vw{entry}{/ˈentri/}{n.}

\sense{n.}{进入；参与，加入；参赛}
\col{entry into sth}
{The thieves gained entry into the building through a back window.}
\col{no entry 禁止入内}
{The sign on the door read: ``No Entry''.}
\col{entry form/level 报名表/入门级别}
{Please complete the entry form and return it by Friday.}

\begin{gram}{entry / entrance / entryway 的"入口"细分}
\textbf{entry}（入口；参赛——通用；多指门/小道：\textit{the entry to the building}）；\textbf{entrance}（入口；进入——多指大堂/正式：\textit{the main entrance}）；\textbf{entryway}（门廊；入口通道——多指美式）。entry 也用作"词典条目"：\textit{a dictionary entry}。
\end{gram}

\begin{origin}
15 世纪由 \textit{enter}（进入）+ \textit{-y} 派生。enter 13 世纪借自古法语 \textit{entrer}，拉丁语 \textit{intrare}（进入），源于 \textit{intra}（在内、内部）。同源词 \textbf{enter}、\textbf{entry}、\textbf{entrance}（入口）、\textbf{enterprise}（企业；\textit{inter-} 间 + \textit{prehendere} 抓$\to$ 抓在中间$\to$ 试图）、\textbf{introduce}（介绍；引入；\textit{intro-} 入 + \textit{ducere} 引$\to$ 引入内）。
\end{origin}

\vw{fee}{/fiː/}{n.}

\sense{n.}{费}
\col{an entry/admission fee}
{The museum charges a small admission fee.}
\col{school/tuition fees 学费}
{School fees have risen sharply over the past decade.}

\sense{n.}{报酬}
\col{a fee for sth}
{She charges a fee for her consulting services.}

\begin{lex}{fee / fare / cost / expense / price / charge 的"费"细分}
\textbf{fee}（费——多指服务费/学费：\textit{an entrance fee}）；\textbf{fare}（车费；船费——交通：\textit{a bus fare}）；\textbf{cost}（成本；费用——通用）；\textbf{expense}（开销——多指商业/差旅：\textit{at company expense}）；\textbf{price}（价格——商品）；\textbf{charge}（收费——多指一次性/服务：\textit{a service charge}）。
\end{lex}

\begin{origin}
古英语 \textit{feoh}（牛；财产；钱），原始日耳曼语 *\textit{fehu}（牛；财富）。同源词 \textbf{fee}、\textbf{fief}（封地——领主赐予佃农的"财产"；古英语 \textit{feoh} 的法语形式）、\textbf{feudal}（封建的——\textit{feodum} 采邑）、德语 \textit{Vieh}（牲畜）。"fee" 字面"牛"——古英语中以牛为财富单位。
\end{origin}

\vw{ban}{/bæn/}{vt. \& n.}

\sense{vt.}{明令禁止，取缔}
\col{ban sth}
{The government has banned the use of single-use plastics in shops.}
\col{ban sb from sth/doing sth}
{She was banned from driving for six months.}

\sense{vt.}{禁止做某事}
\col{ban sb from doing sth}
{The film was banned from being shown in several countries.}

\sense{n.}{禁令}
\col{a ban on sth}
{There is now a total ban on smoking in public places.}
\col{lift/impose a ban 解除/实施禁令}
{The government has imposed a ban on the export of live animals.}

\begin{gram}{ban / forbid / prohibit / bar 的"禁止"细分}
\textbf{ban}（明令禁止——多指官方/法律：\textit{ban smoking}）；\textbf{forbid}（禁止——多指个人/规则：\textit{forbid children to leave}）；\textbf{prohibit}（禁止——正式：\textit{prohibit sth by law}）；\textbf{bar}（阻挡——多指物理/抽象：\textit{bar sb from entry}）。ban 偏官方正式，forbid 偏个人/规则。
\end{gram}

\begin{origin}
古英语 \textit{bannan}（命令、传唤、宣告），原始日耳曼语 *\textit{bannan}（宣告；命令）。同源词 \textbf{ban}、\textbf{banner}（旗帜；横幅；\textit{bannan} 宣告$\to$ 宣告主权之物）、\textbf{bandit}（强盗——与 ban 形近但实不同源，bandit 来自意大利语 \textit{bando} 被流放者）、\textbf{abandon}（放弃；\textit{a-} 朝 + \textit{bannan} 宣告$\to$ 宣告在权力之外$\to$ 放弃）、\textbf{banshee}（报丧女妖——爱尔兰盖尔语 \textit{bean sídhe} 精灵之女）。"ban" 字面"公开宣告"——禁令本是被公开宣告的命令。
\end{origin}

% ---------------- 选收的未标 * 词条 ----------------

\vw{pyramid}{/ˈpɪrəmɪd/}{n.}

\sense{n.}{金字塔}
\col{an Egyptian pyramid}
{The Egyptian pyramids have stood for over four thousand years.}
\col{a glass pyramid}
{Visitors enter the Louvre through I.M. Pei's famous glass pyramid.}

\begin{origin}
16 世纪借自拉丁语 \textit{pyramis}，希腊语 \textit{puramis}（金字塔；小麦蛋糕——外形相似），最终词源不详。同源词 \textbf{pyramid}、\textbf{pyramidal}（金字塔形的；\textit{pyramidal structure} 金字塔结构）。希腊人可能用 \textit{puramis} 既表建筑又表面包——两者皆底宽顶尖。
\end{origin}

\vw{clay}{/kleɪ/}{n.}

\sense{n.}{黏土，陶土}
\col{a clay pot/model}
{She made a small clay pot on the potter's wheel.}
\col{fire/foot of clay 致命弱点}
{Even the greatest leaders have feet of clay.}

\begin{origin}
古英语 \textit{clǣg}（黏土），原始日耳曼语 *\textit{klaig-}（粘）。同源词 \textbf{clay}、\textbf{clayey}（含黏土的）、\textbf{clay pigeon}（飞碟射击的泥鸽——泥制靶子）、德语 \textit{Lehm}（泥）、\textit{clean}（清洁——与 clay 形近但实不同源，clean 来自古英语 \textit{clǣne} 干净）。"feet of clay"（泥足）源自《但以理书》——巨像的脚是"半铁半泥"，寓意伟大者也有致命弱点。
\end{origin}

\vw{cloth}{/klɒθ/}{n.}

\sense{n.}{织物，布料}
\col{a piece of cloth}
{She wiped the table with a piece of cloth.}
\col{clothes/cloth 区别}
{\textit{Cloth} 是不可数名词（\textit{a piece of cloth}）；\textit{clothes} 是复数名词（\textit{new clothes}）；\textit{clothing} 是不可数集合名词（\textit{warm clothing}）。}

\begin{origin}
古英语 \textit{clāþ}（布；衣物），原始日耳曼语 *\textit{klaithaz}。同源词 \textbf{cloth}、\textbf{clothes}、\textbf{clothing}、\textbf{clotheshorse}（晒衣架；注重穿着的人）、\textbf{clad}（穿衣的——古英语 \textit{clǣþan} 穿衣）、德语 \textit{Kleid}（衣服）。
\end{origin}

\vw{shore}{/ʃɔː(r)/}{n.}

\sense{n.}{岸，滨}
\col{on shore}
{The sailors were finally back on shore after three months at sea.}
\col{the shores of ...}
{She grew up on the shores of Lake Michigan.}

\begin{lex}{shore / beach / coast / seaside / bank 的"岸/滨"细分}
\textbf{shore}（岸——海/湖/河的岸：\textit{on the shore}）；\textbf{beach}（海滩——多指沙质：\textit{a sandy beach}）；\textbf{coast}（海岸——地理/整体：\textit{the west coast}）；\textbf{seaside}（海边——休闲语境：\textit{at the seaside}）；\textbf{bank}（岸——多指河岸：\textit{the river bank}）。shore 偏自然，beach 偏休闲，coast 偏地理，bank 偏河流。
\end{lex}

\begin{origin}
14 世纪中古英语 \textit{schore}，可能与中世纪荷兰语 \textit{schor}（沙丘；高地）相关。同源词 \textbf{shore}、\textbf{shore up}（支撑；加固）、德语 \textit{Schor}（陡岸）、荷兰语 \textit{school}（岸）。"shore" 与 "score"（得分）形近但实不同源。
\end{origin}

\vw{tomb}{/tuːm/}{n.}

\sense{n.}{坟墓}
\col{an ancient/empty tomb}
{The ancient tomb was discovered by accident during road construction.}
\col{the tomb of the unknown soldier 无名战士墓}
{She laid a wreath at the tomb of the unknown soldier.}

\begin{origin}
13 世纪借自古法语 \textit{tombe}，晚期拉丁语 \textit{tumba}（坟），希腊语 \textit{tumbos}（土堆；坟）。同源词 \textbf{tomb}、\textbf{tombstone}（墓碑）、\textbf{entomb}（埋葬；\textit{en-} 入 + \textit{tomb}）、\textbf{tumulus}（坟墩——\textit{tumba} $\to$ 土墩）、\textbf{catacomb}（地下墓穴；\textit{cata-} 下 + \textit{tumba} $\to$ 地下土堆）。"tomb" 与 "tummy"（肚子）形近但实不同源——tummy 来自 stomach 的婴儿语形式。
\end{origin}

\vw{wool}{/wʊl/}{n.}

\sense{n.}{毛织物；羊毛}
\col{a ball of wool 毛线球}
{She knitted a sweater from a ball of wool.}
\col{pull the wool over sb's eyes 蒙骗某人}
{Don't let him pull the wool over your eyes — he doesn't really mean what he says.}

\begin{origin}
古英语 \textit{wull}（羊毛），原始日耳曼语 *\textit{wullō}。同源词 \textbf{wool}、\textbf{woollen}（羊毛的；\textit{woollen cloth} 毛料布）、\textbf{woolly}（羊毛的；模糊的；\textit{a woolly mind} 模糊的头脑）、德语 \textit{Wolle}（羊毛）。"pull the wool over one's eyes" 源自 19 世纪流行的男士羊毛假发——扯下假发"蒙住"对方眼睛。
\end{origin}

\vw{cave}{/keɪv/}{n.}

\sense{n.}{山洞，洞穴}
\col{explore a cave}
{The children spent the afternoon exploring caves along the beach.}
\col{caveman 穴居人}
{The painting depicts a caveman hunting with a spear.}

\begin{origin}
13 世纪借自古法语 \textit{cave}，拉丁语 \textit{cava}（中空；洞穴），是 \textit{cavus}（空的）的阴性形式。同源词 \textbf{cave}、\textbf{cavern}（大洞穴；\textit{cava} 增强形式）、\textbf{cavity}（腔；洞；\textit{a cavity in a tooth} 龋齿）、\textbf{concave}（凹的；\textit{con-} 加强 + \textit{cavus}）、\textbf{conceal}（隐藏；\textit{con-} 完全 + \textit{celare} 隐藏——与 cave 形近但实不同源）、\textbf{caveat}（警告；\textit{cavus} 空$\to$ "当心别踩空"）。
\end{origin}

\vw{milestone}{/ˈmaɪlstəʊn/}{n.}

\sense{n.}{重要阶段，里程碑}
\col{a milestone in sth}
{The invention of the printing press was a milestone in the history of communication.}
\col{reach a milestone 达到重要阶段}
{She has just reached a milestone in her career — her 50th novel.}

\begin{lex}{milestone / landmark / watershed / turning point 的"转折"细分}
\textbf{milestone}（里程碑——多指重大成就：\textit{a career milestone}）；\textbf{landmark}（地标；里程碑——可指建筑物：\textit{a landmark ruling} 重要判决）；\textbf{watershed}（分水岭——多指转折点：\textit{a watershed moment}）；\textbf{turning point}（转折点——多指具体时刻）。milestone 强调"标记里程的石头"，watershed 强调"分水岭"。
\end{lex}

\begin{origin}
古英语 \textit{mīl}（英里）+ \textit{stan}（石头）。古罗马人在每英里（约 1618 米）的位置竖立石头作为里程标记。同源词 \textbf{milestone}、\textbf{mile}（英里）、\textbf{landmark}（地标；\textit{land} 地 + \textit{mark} 标记——与 milestone 同类路径）、\textbf{mileage}（英里数；里程）、\textbf{miles}（很远的距离：\textit{miles away}）。
\end{origin}

\vw{equator}{/ɪˈkweɪtə(r)/}{n.}

\sense{n.}{赤道}
\col{on the equator}
{The country lies right on the equator.}
\col{north/south of the equator}
{These tropical regions are located just north of the equator.}

\begin{origin}
14 世纪借自中古拉丁语 \textit{aequator}（平等者；赤道），是拉丁语 \textit{aequare}（使相等）的派生形式，源于 \textit{aequus}（平等的）。同源词 \textbf{equator}、\textbf{equatorial}（赤道的；\textit{equatorial climate} 赤道气候）、\textbf{equal}（平等的；\textit{aequus} 英语化）、\textbf{equality}（平等）、\textbf{equate}（等同）、\textbf{adequate}（充足的；\textit{ad-} 朝 + \textit{aequus} $\to$ 与……平等的$\to$足够的）、\textbf{equilibrium}（平衡；\textit{libra} 天平 + \textit{aequus}）。
\end{origin}

\end{document}
